\documentstyle[%


righttag%


,11pt%


,amssymb%


,russian%               remove in Eng.


,izvru%                 Set izven% in Eng.


]{amsart}





% 


\newcommand{\Alg}{\operatorname{Alg}}


\newcommand{\Conv}{\operatorname{Conv}}


\newcommand{\Id}{\operatorname{Id}}


\newcommand{\f}{\varphi}


\newcommand{\e}{\varepsilon}


\newcommand{\Fr}{\operatorname{Fr}}


\newcommand{\R}{\frak R}


\newcommand{\K}{\frak K}


\newcommand{\M}{\frak M}


\newcommand{\U}{\frak U}


\newcommand{\V}{\frak V}


\newcommand{\E}{\frak E}


\newcommand{\F}{\frak F}


\newcommand{\C}{\frak C}


\newcommand{\I}{\frak I}


\newcommand{\W}{\frak W}


\newcommand{\ep}{\varepsilon}


\newcommand{\sym}{\Sigma}


\newcommand{\s}{\sigma}


\renewcommand{\o}{\omega}


\renewcommand{\t}{\tau}


\newcommand{\al}{\alpha}


\newcommand{\be}{\beta}


\newcommand{\g}{\gamma}


\newcommand{\la}{\lambda}


\renewcommand{\o}{\omega}


\newcommand{\tn}{\otimes}


\newcommand{\A}{\frak A}


\newcommand{\ti}[1]{\widetilde{#1}}


\renewcommand{\P}{\bold K}


\newcommand{\tts}[1]{}





\theoremstyle{plain}


\theorembodyfont{\it}


\newtheorem{lemnonum}{\hskip\parindent Лемма}


\renewcommand{\thelemnonum}{}





%\hoffset=-15mm%                  For Eng.


\setcounter{page}{61}


\god{2002}


\nomer{5~(480)} %                        Remove in Eng.


%\nomer{Vol.\,46, No.\,5}%               Set in Eng.


%\str{\pageref{begin}--\pageref{end}}%  Set in Eng.


\udk{512}





\author{\it С.Н.\,Тронин}


% NB:   ^^ remove in Eng.


\title{Операды в категории конвексоров. II}


\thanks{Работа выполнена при поддержке Российского фонда


фундаментальных исследований (грант \No\,99-01-00469).}





\date{}


%\pagestyle{myheadings}%         Set in Eng.


\pagestyle{plain} %              Remove in Eng.


\begin{document}


%\thispagestyle{plain}%          Set in Eng.


\maketitle


\label{begin}











Данная работа является непосредственным продолжением \cite{T0}.


Результаты обеих частей анонсированы в \cite{T4}. Во второй части 


вводится конвексорный аналог линейных операд: коноперады, т.\,е. операды,


все компоненты которых являются конвексорами, а операции композиции


полилинейны в конвексорном смысле.


Примерами таких операд являются операды многомерных


стохастических и двоякостохастических матриц, первые компоненты которых


--- множества обычных квадратных стохастических и двоякостохастических матриц.


По-видимому, специалисты по теории вероятностей только начинают


искать подходы к изучению этих объектов (см., напр., \cite{Max}).


После подходящей переформулировки вероятностные автоматы становятся элементами


алгебр над операдами многомерных стохастических матриц. 


Основной же результат данной


работы состоит в характеризации многообразий коналгебр над коноперадами:


это в точности те многообразия, которые определяются конвексорными аналогами


полилинейных тождеств. Это аналогично ситуации, имеющей место в


линейном случае \cite{T1}--\cite{T}.








$\Delta$-линейной называется операда $\R$, для которой


соответствие $n\mapsto \R(n)$ является функтором в $\Conv$,


а отображения композиции $\Delta$-полилинейны.


Следуя работам \cite{Sc1}, \cite{Sc2}, такую операду еще можно было бы


назвать коноперадой.


Соответственно, алгеброй над $\Delta$-линейной операдой $\R$ (или,


по аналогии с \cite{Sc1}, \cite{Sc2}, $\R$-коналгеброй) будем называть


$\R$-алгебру $A$, которая является конвексором, в которой все операции


композиции $\R(n)\times A^n\to A$ \ $\Delta$-полилинейны.


Гомоморфизмы таких алгебр будут предполагаться $\Delta$-линейными.


В остальной части работы $\Alg$-$\R$ будет обозначать категорию


алгебр над $\Delta$-линейной операдой.


Категория алгебр и гомоморфизмов $\Alg$-$\R$ является 


многообразием $\Delta$-линейных мультиоператорных алгебр


в смысле \cite{T0}. Сама операда $\R$ рассматривается в качестве


сигнатуры, причем для всех $n\ge 1$ \ $\R(n)$ есть множество символов


$n$-арных $\Delta$-полилинейных операций.





Простейший пример $\Delta$-линейной операды: операда $\C$, у которой


для каждого $n$ множество $\C(n)$ одноэлементно, $\C(n)=\{ \o_n \} $.


Такое множество обладает естественной структурой конвексора


$\al(\o_n,\o_n)=\o_n$. Операция композиции определяется


единственно возможным способом


$\o_m\o_{n_1}\dotsc\o_{n_m}=\o_{n_1+\dotsb+n_m}$.


Ясно, что эти отображения $\Delta$-полилинейны. Для любой подкатегории


$\P$ категории конечных ординалов $\bold{FSet}$, удовлетворяющей условиям,


сформулированным в \cite{T0},


на $\C$ тривиальным образом определяется структура $\P$-операды.








Напомним, что два многообразия алгебр $\frak V_1$ и $\frak V_2$


называются рационально эквивалентными,


если они эквивалентны как категории, причем функтор эквивалентности


сопоставляет алгебре из ${\frak V}_1$, являющейся множеством $A$


с некоторым семейством операций, алгебру из ${\frak V}_2$, являющуюся тем


же самым множеством $A$, но уже, возможно, с другим набором операций.


Гомоморфизмы из ${\frak V}_1$ как отображения множеств переходят


в гомоморфизмы ${\frak V}_2$, являющиеся теми же самыми отображениями.


Обратный функтор действует точно таким же образом. Таким 


образом, алгебры двух многообразий, по сути, это одни и те же множества,


а операции, задающие структуру алгебр двух видов, выражаются друг через друга


(являются взаимно производными). Рациональная эквивалентность многообразий


равносильна изоморфности категорий их свободных


алгебр (достаточно даже свободных


алгебр конечного ранга). Если многообразия $\V_1$, $\V_2$ представлены


каким-то образом в виде подмногообразий многообразия всех $\Omega$-алгебр, то


для их рациональной эквивалентности достаточно, чтобы их свободные алгебры со


счетными базисами были изоморфны как $\Omega$-алгебры.


%% Этот факт будет


%% использован в следующем параграфе.








\begin{lemnonum}


Категория


$\Alg$-$\C_{\sym}$ рационально эквивалентна категории всех коммутативных


ассоциативных конколец $($возможно, без единицы$)$.


Категория $\Alg$-$\C_{\Id}$ рационально эквивалентна категории всех


ассоциативных конколец $($возможно, без единицы$)$.


\end{lemnonum}





\begin{pf}


Определение конколец см. в \cite{Sc2}. Это конвексоры, на которых задана


бинарная ассоциативная операция $x*y$, являющаяся $\Delta$-билинейной.


Лемма доказывается практически так же, как и аналогичное утверждение для


случая, когда конвексорная структура на $\C$ и на алгебрах не


учитывается. Хорошо известно, что в этом случае получаются


соответственно многообразия всех


коммутативных полугрупп (без единицы) и всех полугрупп.


Приведем основные этапы рассуждений.


Положим $\o=\o_2$, $\e=\o_1$.


Если $A$ есть $\C$-алгебра, то определим бинарную операцию


$A\times A\to A$ следующим образом: $x*y=\o (xy)$. Эта операция


$\Delta$-билинейна по определению. Легко проверяется


ассоциативность: $x*(y*z)=$ $\o(\e\o)(xyz)=\o(\o\e)(xyz)=(x*y)*z$.


Если определено действие групп подстановок и $\s=(1\ 2)$ --- транспозиция,


то $\s\o (x,y)=\o(y,x)$, но $\s\o=\o$, откуда $x*y=y*x$.





Обратно, если


дано конкольцо $A$ с операцией умножения $x*y$, то


полагаем $\o_n(x_1,\dotsc,x_n)=x_1*\dotsb *x_n$. Проверка свойств


алгебры над операдой проводится без затруднений. Ясно, что гомоморфизмами


будут и в том и в другом случае одни и те же отображения. В случае, когда


$\P=\Id$, т.\,е. использовать подстановки запрещено, коммутативности не


будет, но ассоциативность умножения $x*y$ сохраняется, и получаем


категорию ассоциативных конколец.


\end{pf}








Приведем еще несколько


более или менее нетривиальных примеров $\Delta$-линейных операд.


Положим $\I(n)=I=[0,1]$ для всех


$n=1,2,\dotsc$, и пусть


$\I(m)\times \I(n_1)\times \dotsb \times \I(n_m)\to \I(n_1+\dotsb +n_m)$


определяется как отображение, сопоставляющее набору чисел из единичного


отрезка $(\la,\la_1,\dots,\la_m)$ их произведение $\la\la_1\dots\la_m$.


Для произвольного $f:[n]\to [m]$


и $\la\in \I(n)=I$ полагаем $f\la=\la\in \I(m)=I$.


Свойства операды проверяются без труда.








Пусть $X$ --- произвольное множество.


Построим операду многомерных стохастических матриц $SM=SM(X)$.


Положим $SM(n)$ равным множеству всех отображений $\o$ из $X\times X^n$


в $[0,1]$ таких, что при любом $\ov{y}\in X^n$ имеет место


$\o(x,\ov{y})=0$ для почти всех $x\in X$ и


$\sum\limits_{x\in X} \o(x,\ov{y})=1$.


Структура конвексора на $SM(n)$ вводится следующим образом:


если $\al\in [0,1]$, $\o_1,\o_2:X\times X^n\to [0,1]$, то


$\al(\o_1,\o_2)(x,x_1,\dotsc,x_n)=$


$\al(\o_1(x,x_1,\dotsc,x_n),\o_1(x,x_1,\dotsc,x_n))$.


Определим действие $\sym_n$


на $SM(n)$, полагая $(\s\o)(x,\ov{y})=\o(x,\ov{y}\s)$. Пусть


$\o_i\in SM(n_i)$, $\o\in SM(m)$, $1\le i \le m$, $\ov{z}_i \in X^{n_i}$,


$x\in X$.


Полагаем


$$


(\o\o_1\dotsc\o_m)(x,\ov{z}_1\dotsc\ov{z}_m)=


\sum_{y_1,\dotsc,y_m\in X}\,\o(x,y_1\dotsc y_m) \o_1(y_1,\ov{z}_1)\dotsc


\o_m(y_m,\ov{z}_m) .


$$





Рассмотрим также множество распределений на $X \times U$, т.\,е.


отображений вида $\gamma : X \times U\to [0,1]$ таких, что


$\gamma(x,u)=0$ для любого $u \in U$ для почти всех


$x\in X$ и $\sum\limits_{x \in X} \gamma(x,u)=1$. Обозначим это множество


через $A_U$ (или $A_{X,U}$). Структура конвексора на


$A_{X,U}$ хорошо известна \cite{Sc2}: $\al(\g_1,\g_2)(x,u)=


\al\g_1(x,u)+(1-\al)\g_2(x,u)$.


Определим композицию $SM$ и $A_U$ следующим образом. Пусть $\o\in SM(n)$,


$\g_1,\dotsc,\g_n\in A_U$. Тогда


$$


\o(\g_1 \dotsc \g_n)(x,u)=\sum \limits_{y_1\dotsc y_n}


\o(x,y_1\dots, y_n)\g_1(y_1,u)\dotsc\g_n(y_n,u).


$$





Прямой проверкой доказывается





\begin{thm}


$SM(X)$ есть $\Delta$-линейная $\sym$-операда.


$A_{U,X}$ есть коналгебра над операдой $SM(X)$.


\end{thm}





Фактически при бесконечном $X$ можно определить еще одну операду такого же


типа, рассматривая те отображения $\o$ из $X\times X^n$


в $\K(n)$, у которых для каждого $x\in X$ имеет место


$\o(x,\ov{y})=0$ для почти всех $\ov{y}\in X^n$, и


$\sum\limits_{\ov{y}\in X^n} \o(x,\ov{y})=1$.


Введенная нами композиция превращается по сути в умножение многомерных


матриц, определенное в \cite{Gas}. ``Бескоординатная'' форма


этой операды, называемая иногда операдой эндоморфизмов \cite{GK}, такова.


Берется конвексор $V$,


$\E_V$ состоит из компонент $\E_V(n)=\Conv(V^{\otimes n}, V)$,


а композиция определяется как


$\o\o_1\o_2\dotsc\o_m=\o(\o_1\otimes\o_2\otimes\dotsb


\otimes \o_m)$. \ $\sym_n$ действует следующим образом:


$\s \o(v_1\otimes\dotsb\otimes v_n)=\o(v_{\s (1)}\otimes\dotsb


\otimes v_{\s (n)})$. Легко убедиться, что для любой


$\Delta$-линейной $\sym$-операды $\R$ задание структуры $\R$-алгебры на


$V$ равносильно заданию гомоморфизма операд $\R\to \E_V$.





С аналогичного примера, рассмотренного в \cite{Art}, судя по всему, и


следует отсчитывать начало изучения (линейных)


операд как самостоятельных объектов. Работа \cite{May}, которую


принято пока считать исходным пунктом развития теории, вышла спустя


более чем два года после \cite{Art}.





Определим еще операду двоякостохастических многомерных матриц


$DSM=DSM(X)$, полагая $DSM(n)$ состоящим из всех тех


$\o\in SM(n)$, $\o:X\times X^n\to [0,1]$, для которых


при любом $x\in X$ имеет место равенство


$\sum\limits_{\ov{y} \in X^n} \o(x,\ov{y})=1$.





\begin{thm}


$DSM(X)$ есть $\Delta$-линейная $\sym$-операда.


\end{thm}





{\bf Доказательство} проводится проверкой определения.$\quad\square$


\medskip





Пусть $B=\{ 0,1 \}$ --- двухэлементная булева алгебра.


Зафиксируем снова множество $X$ и определим операду булевских стохастических


матриц $SBM=SBM(X)$. Положим


$SBM(n)$, равным множеству всех отображений вида $\o :X \times X^n \to B$


таких, что при любом $ \ov{y} \in X^n$


будет $\o (x,\ov{y})=0$ для почти всех $ x\in X$


и $\sum \limits_{x \in X} \o (x,\ov{y})=


\mathop{\max}\limits_{x \in X} \o(x,\ov{y})=1 $.


Определим на $SBM$ структуру операды. Положим


\begin{gather*}


SBM(m)\times SBM(n_1)\times \dotsb \times SBM(n_m) \to


SBM(n_1+\dotsb + n_m),\\


(\o \o _1 \dotsc \o _m)(x,\ov{z}_1\dotsc \ov{z}_m)=


\sum_{y_1 \dotsc y_m \in X}


\o(x,y_1 \dotsc y_m)\o _1(y_1,\ov{z}_1)\dotsc \o_m(y_m,\ov{ x}_m).


\end{gather*}


Здесь сумма означает взятие максимума, а


произведение --- минимума.





Напомним, что каждая верхняя полурешетка обладает структурой конвексора:


достаточно положить $\al(x,y)=x+y$ при $\al\in (0,1)$, $1(x,y)=x$,


$0(x,y)=y$. Отсюда следует, что все $SBM(n)$ являются конвексорами,


и определенная только что операция композиции $\Delta$-полилинейна.


Определим $ BA=BA_{U,X}$ как множество всех отображений


$\g: X\times U\to B$ таких, что для любого $ u\in U$ выполняется $\g(u,x)=0$


для почти всех $ x\in X$ и


$\sum \limits_{x \in X} \g (u,x)=1$.


Композиция $SBM(n)\times BA^n\to BA$ определяется так же, как и выше для


$SM$ и $A$. Зададим отображение


$Q:[0,1] \to \{0,1\}$, полагая $Q(0)=0$, $Q(a)=1$ при $a>0$.


Построим семейство отображений


$Q_n : SM(n)\to SBM(n)$ следующим образом. Если


$A \in SM(n)$, то $Q_n(A)\in SBM(n)$ определяется в виде


$Q_n(A)(y,x_1\dotsc x_m)= Q(A(y,x_1\dotsc x_m))$.


Следуя (\cite{Buh}, c.\,32),


можно назвать $Q_n(A)$ булевскими шаблонами многомерных стохастических матриц.


Там, где это не приводит к


недоразумениям, будем обозначать через


$Q$ все семейство $Q_n$, $n=1,2,\dotsc$





\begin{thm}\label{m0}


Семейство $SBM$ с определенной выше композицией есть операда.


На ней также существует естественная структура


$\Delta$-линейной операды.


$BA_{U,X}$ есть алгебра над $SBM(X)$. Семейство $Q$


есть гомоморфизм $\Delta$-линейных операд,


$Q:SM \to SBM$.


\end{thm}





\begin{pf}


Все утверждения, кроме последнего, доказываются непосредственной


проверкой определений.


$\Delta$-линейность отображения $Q_n: SM(n)\to SBM(n)$ также очевидна.


Проверим, что $Q(AA_1\dotsc A_m)=Q(A)Q(A_1)\dotsc Q(A_m)$.


Пусть $\ov{z}_i\in X^{n_i}$, $1\le i\le m$. Тогда


$Q(AA_1\dotsc A_m)(x,\ov{z}_1\dots\ov{z}_m)=0$ в том и только том случае,


когда $(AA_1\dotsc A_m)(x,\ov{z}_1\dots\ov{z}_m)=0$. Это равносильно


тому, что для любого семейства $y_1,\dotsc, y_m\in X$ по крайней мере одно


из чисел $A(x,y_1\dotsc y_m), A_1(y_1,\ov{z}_1), \dotsc,


A_m(y_m,\ov{z}_m)$ равно нулю, или что по крайней мере одна из


компонент булевских шаблонов $Q(A)(x,y_1\dotsc y_m), Q(A_1)(y_1,\ov{z}_1),


\dotsc, Q(A_m)(y_m,\ov{z}_m)$ является нулевой. С другой стороны,


\begin{multline*}


Q(A)Q(A_1)\dotsc Q(A_m)(x,\ov{z}_1\dotsc\ov{z}_m)= \\


=\max_{y_1,\dotsc,y_m}


(\min(Q(A)(x,y_1\dotsc y_m), Q(A_1)(y_1,\ov{z}_1),


\dotsc,Q(A_m)(y_m,\ov{z}_m))),


\end{multline*}


и это выражение равно нулю в том и только том случае, если для


любого семейства $y_1,\dotsc, y_m{\in} X$ %%%%


$$


\min(Q(A)(x,y_1\dotsc y_m), Q(A_1)(y_1,\ov{z}_1),\dotsc,


Q(A_m)(y_m,\ov{z}_m))=0,


$$


т.\,е. равно нулю хотя бы одно из


$Q(A)(x,y_1\dotsc y_m), Q(A_1)(y_1,\ov{z}_1), \dotsc,


Q(A_m)(y_m,\ov{z}_m)$.


Итак, обе сравниваемые функции принимают нулевые значения на одних и тех же


значениях аргумента. Следовательно, на остальных значениях аргумента они


обязаны быть равными единице.


\end{pf}





Последнее утверждение теоремы 3 --- это


многомерный аналог леммы 1.3.6 из книги (\cite{Buh}, с.\,33).


Легко заметить, что аналогичный факт справедлив и для алгебр распределений:


для них также можно определить булевские шаблоны, причем соответствующие


отображения оказываются гомоморфизмами алгебр над операдами.


Напомним определение вероятностного автомата \cite{Buh}. Пусть


$X$ --- множество входных сигналов, $Y$ --- множество


выходных сигналов, $Q$ --- множество состояний. Вероятностный автомат


--- это отображение


$\g : X\times Q\times Y\times Q \to [0,1]$ такое, что


для каждой фиксированной пары $(x,q)\in X\times Q$


имеют место равенства $\g(x,q,y,q')=0$ для почти всех


$(y,q')\in Y\times Q$ и $\sum\limits_{(y,q) \in Y\times Q}


\g(x,q,y,q')=1$. Подразумевается, что


$\g(x,q,y,q')$ есть условная вероятность перехода из события $q$


при входном сигнале $x$ в состояние $q'$ при выходном сигнале $y$.


Обозначим через $A=A_{X,Q,Y}$ множество всех таких автоматов,


$Z=Y\times Q$, $SM=SM(Z)$. Oпределим отображения композиции


$SM(n)\times A^n\to A$


следующим образом:


$$


(\o \g_1 \dotsc \g_n)(x,q,y,q')=


\sum \begin{Sb} y'_1 \dotsc y'_n \\  q'_1\dotsc q'_n\end{Sb}


\o((y,q'),(y'_1,q'_1)\dotsc (y'_n,q'_n))\g_1(x,q,y'_1,q'_1)\dotsc


\g_n(x,q,y'_n,q'_n).


$$





\begin{thm}


Определенные выше операции композиции превращают


$A_{X,Q,Y}$ в коналгебру над коноперадой $SM$.


\end{thm}





{\bf Доказательство} проводится прямой проверкой определения.$\quad\square$


\medskip





Отметим, что сама операда $\Delta$ не является $\Delta$-линейной.








\begin{thm}\label{m1}


Свободная $\Delta$-линейная алгебра с базисом $X$ 


в многообразии $\Alg$-$\R_{\Id}$


имеет вид


$\Fr_{\R}(X)=\coprod\limits_{n=1}^{\infty} \R(n)\otimes T^n(X)$,


где $T^n(X)$ --- свободный конвексор с базисом $X^n$, а


$\coprod$ --- копроизведение в категории $\Conv$.


Если $\P$ --- подкатегория категории конечных множеств,


обладающая свойствами, сформулированными в {\em ([1], \S\,1)}, и $\R$ есть


$\P$-операда, то


в многообразии $Alg$-$\R_{\bold K}$, которое является подмногообразием


$\Alg$-$\R_{\Id}$, свободная алгебра $\Fr_{\R,{\bold K}}(X)$


есть факторалгебра $\Fr_{\R}(X)$ по конгруэнции,


порожденной всеми парами вида


$((f\o)\otimes (x_1,\dotsc, x_m)) \sim


(\o \otimes (x_{f(1)}, \dotsc,x_{f(n)}))$,


где $f\in \P(n,m)$,


$\o\in \R(n)$, $x_1,\dotsc, x_m\in X$.


В случае $\P=\sym$


свободные алгебры в $Alg$-$\R_{\sym}$ имеют вид


$\Fr_{\R}(X)=\coprod\limits_{n=1}^{\infty} \R(n)\otimes_{\sym_n} T^n(X)$.


\end{thm}





\begin{pf}


Для линейных $\sym$-операд этот факт хорошо известен


(см. \cite{T}, \cite{GK}).


В $\Delta$-линейном случае все получается точно так же. Напомним лишь,


как действует операда $\R$ на алгебре $\Fr_{\R}(X)$:


$$


\o(\o_1\otimes \ov{x}_1, \dotsc, \o_m\otimes \ov{x}_m)=


(\o\o_1\dotsc \o_m)\otimes (\ov{x}_1\dotsc \ov{x}_m).


$$


Здесь


$\o\in \R(m)$, $\o_i\in\R(n_i)$, $\ov{x}_i\in X^{n_i}$, $1\le i \le m$.


В случае произвольной категории $\P$ все следует из определения.


\end{pf}





Далее рассматриваются только $\sym$-операды, и это не будет особо


оговариваться. Определим свободную $\Delta$-линейную операду


с базисом $\Omega=\mathop{\cup}\limits_{n=1}^{\infty} \Omega_n$


как операду $\F=\F(\Omega)$ вместе с вложениями


$\Omega_n\subseteq \F(n)$ такими, что для любой $\Delta$-линейной операды


$\R$ любое семействo отображений вида $\Omega_n\to \R(n)$, $n=1,2,\dotsc$,


единственным способом продолжается до гомоморфизма операд $\F\to \R$.


Свободные операды существуют и строятся примерно так же, как и абсолютно


свободные алгебры. А именно, сначала индуктивно строится нелинейная


свободная операда $\W$. Основание индукции: $\Omega_n \subset \W(n)$


для всех $n\ge 1$, и, кроме того, постулируется существование


единицы (``пустого слова'') $\ep\in \W(1)$. Далее, если уже построен


терм (слово) $t\in \W(n)$, то определен терм $\s t\in \W(n)$ для любого


$\s\in \sym_n$, а если определены термы $w\in \W(m)$, $w_i \in \W(n_i)$,


то по ним определяется терм $ww_1\dotsc w_m\in \W(n_1+\dotsb+n_m)$.


Чтобы получилась операда, необходимо профакторизовать множество термов по


наименьшему отношению эквивалентности, содержащему все тождества


из определения операды в \cite{T0}.


Стандартным образом показывается выполнимость


универсального свойства. $\Delta$-линейная операда $\frak F$ есть семейство


свободных конвексоров с базисами $\W(n)$. Операции композиции определяются


``по $\Delta$-линейности''.








\begin{thm}


Многообразие


$\Alg$-$\F$ рационально эквивалентно многообразию всех


$\Delta$-линейных $\Omega$-алгебр, так что $\Fr_{\F}(Z)\cong \Fr_{\Omega}(Z)$


для всех базисов $Z$.


Компонента $\F (n)$ свободной $\Delta$-линейной $\sym$-операды


с базисом $\Omega$ изоморфна $($как конвексор и как $\sym_n$-множество$)$


подмножеству всех однородных $n$-полилинейных элементов абсолютно свободной


$\Delta$-линейной $\Omega$-алгебры $\Fr_{\Omega}(X)$, где 


$X=\{x_1,\dotsc, x_n\}$.


\end{thm}





\begin{pf}


Рассмотрим $\F$-алгебру $A$. Так как по построению


$\Omega_n\subset \F(n)$, то из наличия $\Delta$-полилинейных композиций


$\F(n)\times A^n\to A$ следует, что для каждого $\o\in \Omega_n$ определено


$\Delta$-полилинейное отображение $\o^A: A^n\to A$, и, таким образом,


каждая $\F$-алгебра естественным образом превращается в $\Omega$-алгебру.


Обратно, задание структуры $\Delta$-линейной $\Omega$-алгебры на $A$


эквивалентно заданию семейства отображений вида $\Omega_n\to \E_A(n)$ для всех


$n=1,2,\dotsc$ Но согласно определению свободной операды


по этому семейству отображений единственным образом строится гомоморфизм


операд $\F\to \E_A$, что равносильно заданию на $A$ структуры $\F$-алгебры.


Легко проверяется, что эти соответствия структур алгебр взаимно обратны


и что в обоих случаях гомоморфизмами являются одни и те же


отображения. Таким образом, имеет место изоморфизм категорий именно того


вида, который фигурирует в определении рациональной эквивалентности.





Чтобы установить последнее утверждение, рассмотрим подмножества


$\F'(n)\subset \Fr_{\Omega}(x_1,\dotsc,x_n)$, состоящие из


полилинейных элементов. Так как согласно предыдущему


$\Fr_{\Omega}(x_1,\dotsc,x_n)$ можно отождествить с


$\coprod\limits_{n=1}^{\infty} \F(n)\otimes_{\sym_n} T^n(X)$,


то элементы $\F'(n)$ имеют вид


$\sum\limits_{\s} \al_s f_s\otimes (x_{\s(1)},\dotsc,x_{\s(n)})$,


где $f_s\in \F(n)$, суммирование ведется по некоторому множеству подстановок


$\s\in \sym_n$ и $\sum\limits_{\s} \al_s=1$. Используя


свойства тензорного произведения (аналогичные тем, которые справедливы для


линейного случая), приводим это выражение к виду


$\Big(\sum\limits_{\s} \al_s (\s^{-1}f_s)\Big)\otimes


(x_{1},\dotsc,x_{n})$ или $f\otimes (x_{1},\dotsc,x_{n})$, где


$f\in \F(n)$. Этот элемент можно отождествить с $f$, что дает


искомый изоморфизм. Более формальное обоснование: $T^n({x_1,\dotsc, x_n})$


разлагается в копроизведение циклических конмодулей над конкольцом ---


аналогом групповой алгебры группы $\sym_n$, причем прямые слагаемые


соответствуют орбитам действия $\sym_n$ на $X^n$. В частности, одно


из прямых слагаемых --- свободный конмодуль с базисом $(x_1,\dotsc,x_n)$.


Поскольку функтор тензорного произведения перестановочен с копроизведениями


(ввиду наличия сопряженного, см. [1]) и тензорное произведение


$\F(n)$ над $\sym_n$ на свободный $\sym_n$-конмодуль с базисом из одного


элемента изоморфно $\F(n)$, то указанное выше соответствие


$f \to f\otimes (x_1,\dotsc, x_n)$ является изоморфизмом конвексоров и


даже $\sym_n$-эквивариантным.


$\F'$ превращается в операду следующим образом. Если даны


полилинейные элементы $\o=f\otimes (x_1,\dotsc,x_m)$,


$\o_i=f_i\otimes (x_1,\dotsc, x_{n_i})$, $1\le i \le m$, то


$\o\o_1\dotsc\o_m$ есть результат подстановки в $\o$ вместо $x_i$


элемента $\o_i$ с одновременной перенумерацией переменных в $\o_i$ для всех


$i\ge 2$


$$


(f f_1\dotsc f_m)\otimes (x_1,\dotsc,x_{n_1},x_{n_1+1},\dotsc,x_{n_1+n_2},


\dotsc,x_{n_1+\dotsb+n_{m-1}+1},\dotsc, x_{n_1+\dotsb+n_m}).


$$


Используя данное при доказательстве теоремы 5 описание


композиции в $\F$-алгебре $\Fr_{\F}(x_1,\dotsc,x_n)$,


легко убедиться, что изоморфизм конвексоров $\F(n)$ и $\F'(n)$ ---


это изоморфизм операд. Заметим также, что


аналогичным образом можно отождествить $\R(n)$ с подмножеством


формальных полилинейных элементов $\Fr_{\R}(x_1,\dotsc,x_n)$ для любой


операды $\R$.


\end{pf}





Элементы $f\otimes(x_1,\dotsc,x_n)$ будем записывaть в виде


$fx_1\dotsc x_m$. Прежде чем доказывать следующую теорему, напомним,


что под конгруэнцией в алгебре $A$ подразумевается подалгебра


$A\times A$, которая является отношением эквивалентности на $A$.


Если $A$ --- свободная алгебра, то вербальная конгруэнция


на $A$ --- это вполне инвариантная конгруэнция, т.\,е.


для любого эндоморфизма $\varphi:A\to A$ эндоморфизм


$\varphi\times \varphi$ переводит эту подалгебру в себя.


Когруэнцией в операде $\R$ называется подоперада $\U$ в


$\R\times \R$ такая, что все $\U(n)$ являются отношениями эквивалентности


в $\R(n)$. Напомним, что $(\R\times\R)(n)=\R(n)\times\R(n)$, композиция в


операде и действие групп подстановок определяются покомпонентно.








\begin{thm}


Существует изоморфизм между решеткой конгруэнций свободной


$\Delta$-линейной операды $\F$ с базисом $\Omega$ и решеткой


вербальных конгруэнций в абсолютно


свободной $\Delta$-линейной $\Omega$-алгебре $\Fr(X)$


со счетным базисом $X$, порожденных полилинейными элементами. Если $\U$ ---


конгруэнция в $\F$, то соответствующая вербальная конгруэнция $\ti{\U}$


имеет следующий вид\/$:$


$\coprod\limits_{n=1}^{\infty} \U(n)\otimes_{\sym_n} T^n(X)$.


Точнее, имеется в виду образ при отображении


$$


\coprod_{n=1}^{\infty} \U(n)\otimes_{\sym_n} T^n(X)


\rightarrow


\coprod_{n=1}^{\infty}


(\F(n)\times \F(n))\otimes_{\sym_n} T^n(X)


\rightarrow


\bigg(\coprod_{n=1}^{\infty} \F(n)\otimes_{\sym_n} T^n(X)\bigg)^2.


$$


Первое из участвующих здесь отображений индуцировано вложениями


$\U(n)\subseteq \F(n)\times \F(n)$, а второе на порождающих элементах имеет


вид $(\o_1,\o_2)\otimes\ov{x}\mapsto(\o_1\otimes\ov{x},\o_2\otimes \ov{x})$.


Имеет место изоморфизм $\Delta$-линейных $\Omega$-алгебр\/$:$


$\Fr_{\F}(X)/\ti{\U}\cong \Fr_{\F/\U}(X)$.


\end{thm}





\begin{pf}


Пусть дана вербальная конгруэнция $U\subseteq \Fr(X)\times \Fr(X)$,


где $X$ счетно. Рассмотрим множество $P$ всех пар полилинейных


элементов вида $(\o\ov{x}, \mu\ov{x})$, где в $\ov{x}=x_{i_1}\dotsc x_{i_n}$


все $x_{i_1},\dotsc, x_{i_n}$ различны, $\o,\mu\in \F(n)$.


Утверждается, что тогда семейство множеств $\widehat{U}=\U$,


$\U=\{\U(n)\mid n=1,2, \dotsc\}$,


$\U(n)=\{\ (\o,\mu)\in \F(n)^2 \mid


(\o\ov{x}, \mu\ov{x})\in P$ для некоторого $\ov{x} \}$


образует конгруэнцию в операде $\F$. Во-первых, заметим, что


выбор $\ov{x}$ для пары $(\o, \mu)$ произволен в том смысле, что


вместо $\ov{x}=x_{i_1}\dotsc x_{i_n}$ можно взять любое другое


слово $\ov{x}'=x_{j_1}\dotsc x_{j_n}$ с условием, что все


$x_{j_1},\dotsc, x_{j_n}$ --- различные элементы $X$. Это следует из


инвариантности конгруэнции $U$: соответствие $x_{i_k}\mapsto x_{j_k}$,


$1\le k\le n$, продолжается до эндоморфизма $\Fr(X)$, затем до


эндоморфизма $\Fr(X)\times \Fr(X)$, отображающего $U$ в $U$, а пару


$(\o\ov{x}, \mu\ov{x})$ --- в пару $(\o\ov{x}', \mu\ov{x}')$.


Очевидно, что все $\U(n)$ будут отношениями эквивалентности и конвексорами.


Пусть $(\o,\mu)\in \U(m)$, $(\o_i,\mu_i)\in\U(n_i)$, $1\le i\le m$.


Покажем, что $(\o\o_1\dotsc\o_m, \mu\mu_1\dotsc\mu_m)\in$


$\U(n_1+\dotsb+n_m)$.


Выберем $\ov{x}_1,\dotsc,\ov{x}_m$ --- слова в алфавите $X$


так, что $(\o_i\ov{x}_i,\mu_i\ov{x}_i)\in P$, $1\le i\le m$, и в


$\ov{x}_i$, $\ov{x}_j$ нет общих символов для всех $i\ne j$. Это


можно сделать ввиду счетности $X$. Поскольку


$U\subseteq \Fr(X)\times \Fr(X)$ --- подалгебра $\F$-алгебры $\Fr(X)$,


то для любого $\o\in\F(n)$ и любых


$(\o_i\ov{x}_i,\mu_i\ov{x}_i)\in P$, $1\le i\le m$, получим элемент из


$U$


$$


\o(\o_1\ov{x}_1,\mu_1\ov{x}_1)\dotsc (\o_m\ov{x}_m,\mu_m\ov{x}_m)=


((\o\o_1\dotsc\o_m)(\ov{x}_1 \dotsc \ov{x}_m),


(\o\mu_1\dotsc\mu_m)(\ov{x}_1\dotsc\ov{x}_m)).


$$


Но ввиду выбора $\ov{x}_i$ этот


элемент должен принадлежать $P$. Таким образом,


$$


(\o\o_1\dotsc\o_m,\o\mu_1\dotsc\mu_m)\in \U(n_1+\dotsb+n_m).


$$


Пусть $\ov{x}=x_{t_1}\dotsc x_{t_m}$ --- слово в алфавите $X$


такое, что все $x_{t_j}$ различны и $\ov{x}$ не имеет общих символов ни с


одним из $\ov{x}_i$, $1\le i \le m$. Тогда


$(\o\ov{x},\mu\ov{x})\in P\subset U$. Рассмотрим некоторый эндоморфизм


$\Fr(X)$, отображающий $x_{t_j}$ в $\mu_j\ov{x}_j$, $1\le j \le m$.


Тогда $(\o\ov{x},\mu\ov{x}) $ отображается в элемент~$U$


$$


(\o(\mu_1\ov{x}_1\dotsc \mu_m\ov{x}_m),


\mu(\mu_1\ov{x}_1\dotsc\mu_m\ov{x}_m))=


((\o\mu_1\dotsc\mu_m)(\ov{x}_1 \dotsc \ov{x}_m),


(\mu\mu_1\dotsc\mu_m)(\ov{x}_1\dotsc\ov{x}_m)).


$$


Очевидно, однако, что вновь получен элемент из $P$, так что


$(\o\mu_1\dotsc\mu_m,\mu\mu_1\dotsc\mu_m)\in


\U(n_1+\dotsb+n_m)$.


Поскольку $\U(n_1+\dotsb+n_m)$ есть отношение эквивалентности, получаем


требуемое включение


$(\o\o_1\dotsc\o_m, \mu\mu_1\dotsc\mu_m)\in\U(n_1+\dotsb+n_m)$.


Из инвариантности конгруэнции $U$ также непосредственно следует, что $\U(n)$


замкнуты относительно действия групп $\sym_n$.





Из построения соответствия $U\mapsto \widehat{U}=\U$ ясно, что


оно сохраняет включения и произвольные пересечения. Обратное соответствие


$\U\mapsto \ti{U}$ определено в формулировке теоремы. Нетрудная проверка


показывает, что $\ti{U}$ есть конгруэнция.


Проверим ее полную инвариантность.


Пусть гомоморфизм $\varphi: \Fr(X)\to \Fr(X)$ отображает


$x_i$ в $\sum\limits_j \al_{i\, j} (w_{i\, j} \ov{x}_{i\, j})$,


где $w_{i\, j}\in \F$, $0\le \al_{i\, j}\le 1$,


$\sum\limits_j \al_{i\, j}=1$, и


$(u_1 x_1\dotsc x_m, u_2 x_1\dotsc x_m)$ --- порождающий элемент $\ti{U}$


(здесь можно даже не предполагать, что все $x_1,\dotsc, x_m$ различны).


Тогда применение $\varphi$ к этому элементу дает


$(u_1 \f (x_1)\dotsc \f(x_m), u_2 \f (x_1)\dotsc \f(x_m))$,


что преобразуется к виду (с учетом того, как устроено произведение


конвексоров)


$$


\sum_{j_1}\dotsc\mathop{\sum}\limits_{j_m}


\al_{1\, j_1} \dotsc \al_{m\, j_m}


((u_1 w_{1\, j_1}\dotsc w_{m\, j_m})\ov{x}_{j_1\dots j_m},


(u_2 w_{1\, j_1}\dotsc w_{m\, j_m})\ov{x}_{j_1\dots j_m}),


$$


где $\ov{x}_{j_1\dotsc j_m}=\ov{x}_{1\, j_1}\dotsc \ov{x}_{m\, j_m}$.


Остается заметить, что в операде $\F\times \F$


$$


(u_1 w_{1\, j_1}\dotsc w_{m\, j_m},u_2 w_{1\, j_1}\dotsc w_{m\, j_m})=


(u_1,u_2)(w_{1\, j_1}, w_{1\, j_1})\dotsc (w_{m\, j_m},w_{m\, j_m}),


$$


причем $(u_1,u_2)\in \U(m)$ по выбору, а любая пара $(w,w)\in \U(k)$,


т.\,к. все $\U(k)$ --- отношения эквивалентности. Следовательно,


результат записанной выше композиции также принадлежит подопераде


$\U$. Это доказывает вербальность конгруэнции $\ti{U}$.





Покажем взаимную обратность соответствий $U\mapsto \widehat{U}$


и $\U\to\ti{\U}$. В одну сторону это очевидно: полилинейные


элементы из $\ti{\U}$ приводят вновь к операде $\U$.


Проверим, что если $\U=\widehat{U}$, то $U=\ti{\U}$.


До сих пор не использовалось то, что $U$ порождается множеством


$P$ как вербальная конгруэнция. Это означает, что произвольный элемент


$(z_1,z_2)\in U$ представляется в виде линейной комбинации (с


коэффициентами из отрезка $[0,1]$, сумма которых равна единице)


элементов вида $\o(p'_1,p''_1)\dotsc(p'_m,p''_m)$, где $\o\in \F(m)$,


а каждая пара $(p'_i,p''_i)$ получается из некоторого элемента $P$


подстановкой вместо переменных некоторых элементов $\Fr(X)$, причем


для каждого отдельно взятого $i$, $1\le i\le m$, вместо одинаковых


переменных подставляются одинаковые элементы.


С помощью рассуждений, сходных с теми, которые выше были использованы для


доказательства вербальности $\ti{\U}$, заключаем, что каждый элемент


$(p'_i,p''_i)$ представляется в виде линейной комбинации элементов


вида $(q'_i \ov{x}_i,q''_i \ov{x}_i)$, где $(q'_i,q''_i)\in \U=\widehat{U}$.


Следовательно, произвольный элемент $U$ есть линейная комбинация


элементов вида


\begin{equation}


\o (q'_1 \ov{x}_1,q''_1 \ov{x}_1)\dotsc (q'_m \ov{x}_m,q''_m\ov{x}_m)=


((\o q'_1\dotsc q'_m)\ov{x}_1\dotsc\ov{x}_m,


(\o q''_1\dotsc q''_m)\ov{x}_1\dotsc\ov{x}_m). \tag{$\ast$}


\end{equation}


Пара $(\o q'_1\dotsc q'_m,\o q''_1\dotsc q''_m)$ принадлежит $\U$.


Чтобы убедиться в этом, выберем


для каждой пары


$(q'_i,q''_i)$ слово $\ov{y}_i$ в алфавите $X$ таким образом, что


$(q'_i \ov{y}_i,q''_i \ov{y}_i)\in P$, причем


в $\ov{y}_i$ и $\ov{y}_j$ нет общих символов при всех $i\neq j$.


Поскольку $U$ --- подалгебра, то композиция


$$


\o (q'_1 \ov{y}_1,q''_1 \ov{y}_1)\dotsc (q'_m \ov{y}_m,q''_m\ov{y}_m)=


((\o q'_1\dotsc q'_m)\ov{y}_1\dotsc\ov{y}_m,


(\o q''_1\dotsc q''_m)\ov{y}_1\dotsc\ov{y}_m)


$$


принадлежит $U$ и очевидно, что она принадлежит и $P$.


Возвращаясь к элементу $(\ast)$,


заключаем, что он принадлежит $\ti{U}$. Включение $\ti{U}\subseteq U$


очевидно.





Таким образом, получено взаимно однозначное соответствие, сохраняющее


порядок, между двумя решетками, причем отображение $U\mapsto \widehat{U}$


сохраняет произвольные пересечения.


Но так как точная верхняя грань элементов


$U_1$ и $U_2$ в рассматриваемых решетках есть пересечение всех~$U$


таких, что $U_1,U_2\subseteq U$, то имеет место изоморфизм решеток.





Рассмотрим диаграмму


$\U(n)\overset{v}{\subseteq} \F(n)\times \F(n)


\mathop{\overset{\scriptstyle \pi_1}{\longrightarrow}}


\limits_{\overset{\displaystyle \longrightarrow}{\scriptstyle \pi_2}} \F(n)


\overset{c}{\longrightarrow} \F(n)/\U(n)$,


в которой $\pi_1$, $\pi_2$ --- естественные проекции на первый и на второй


сомножители, $v$ --- вложение, а $c$ есть коуравнитель пары


$(\pi_1 v, \pi_2 v)$. Как было отмечено в [1], тензорные произведения


сохраняют коуравнители, так что в следующей диаграмме


$c\otimes 1$ есть коуравнитель пары


$((\pi_1\otimes 1)(v\otimes 1),(\pi_2\otimes 1)(v\otimes 1))$:


$$


\U(n)\underset{\sym_n}{\otimes}


T^n(X)@>{v\otimes 1}>>


(\F(n)\times \F(n))\underset{\sym_n}{\otimes}T^n(X)


\mathop{\overset{\scriptstyle \pi_1\otimes 1}{\longrightarrow}}


\limits_{\overset{\displaystyle \longrightarrow}{\scriptstyle\pi_2\otimes1}}


\F(n)@>{c\otimes 1}>>


((\F(n)/\U(n))\underset{\sym_n}{\otimes} T^n(X)).


$$


Рассмотрим отображение


$d:(\F(n)\times \F(n))\otimes_{\sym_n} T^n(X)


\rightarrow (\F(n)\otimes_{\sym_n} T^n(X))^2$ такое, что


$d((\o_1,\o_2)\otimes \ov{x})=(\o_1\otimes \ov{x},\o_2\otimes \ov{x})$.


Если $p_i: (\F(n)\otimes_{\sym_n} T^n(X))^2 \to \F(n)\otimes_{\sym_n} T^n(X)$,


$i=1,2$, --- естественные проекции, то имеют место очевидные равенства


$p_i d=\pi_i\otimes 1$, $i=1,2$. Отсюда следует, что


$\F(n)/\U(n)\otimes_{\sym_n} T^n(X)$ --- коуравнитель пары


$(p_1 d (v\otimes 1),p_2 d (v\otimes 1))$, т.\,е. фактормножество


$(\F(n)\otimes_{\sym_n} T^n(X))^2$ по конгруэнции, являющейся образом


$d (v\otimes 1)$. Остается заметить, что в любой категории, где


соответствующие операции определены,


взятие копроизведения произвольного семейства диаграмм коуравнителей


есть также диаграмма коуравнителя. Это значит, что свободная алгебра


многообразия $\Alg$-$(\F/\U)$, имеющая вид


$\coprod\limits_{n=1}^{\infty} (\F(n)/\U(n))\otimes_{\sym_n} T^n(X)$,


изоморфна факторалгебре $\Fr_{\F}(X)$ по конгруэнции $\ti{\U}$.


\end{pf}








\begin{thm}


Многообразия $\Delta$-линейных мультиоператорных алгебр,


определяемые $\Delta$-поли\-ли\-ней\-ными тождествами $($и только они$)$


имеют вид $\Alg$-$\R$ для подходящей $\Delta$-линейной операды $\R$.


\end{thm}





\begin{pf}


Пусть дана коноперада $\R$. Представим ее в виде фактороперады


свободной операды $\F$ по конгруэнции $\U$. Тогда $\Alg$-$\R$ можно


вложить в качестве подмногообразия в многообразие $\Alg$-$\F$, которое


можно считать многообразием всех $\Delta$-линейных $\Omega$-алгебр


(выбор $\Omega$ достаточно произволен, это множество должно лишь находится


во взаимно однозначном соответствии с некоторым множеством образующих $\R$).


Из предыдущей теоремы следует, что для счетного $X$ имеет место изоморфизм


между $\Fr_{\R}(X)$ и факторалгеброй $\Fr_{\F}(X)$ по конгруэнции $\ti{\U}$.


Но конгруэнция $\ti{\U}$ порождается полилинейными тождествами, а это 


и означает, что $\Alg$-$\R$ определяется полилинейными тождествами.





Обратно, пусть дано многообразие $\Delta$-линейных $\Omega$-алгебр $\V$,


определяемое полилинейными тождествами, и $F$ есть свободная


алгебра из $\V$ со счетным базисом $X$. Согласно условию $F$ изоморфна


факторалгебре абсолютно свободной алгебры $\Fr_{\F}(X)$ по вербальной


конгруэнции, порожденной полилинейными элементами. Но уже известно,


что эта конгруэнция имеет вид $\ti{\U}$. Из последнего утверждения предыдущей


теоремы следует, что


$\Fr_{\F}(X)/\ti{\U}\cong \Fr_{\F/\U}(X)$ как


$\Delta$-линейные $\Omega$-алгебры. Но это означает, что


многообразия $\V$ и $\Alg$-$(\F/\K)$


совпадают как подмногообразия $\Alg$-$\F$.


\end{pf}








Заметим, что сходным образом автором было ранее доказано


(\cite{T1}--\cite{T}), что для


многообразий линейных мультиоператорных алгебр (над коммутативными кольцами)


класс многообразий, задаваемых полилинейными тождествами, совпадает с


классом многообразий вида $\Alg$-$\R$ (под совпадением понимается, конечно,


рациональная эквивалентность).


Недавно аналогичный факт установлен и для


многообразий супералгебр \cite{T5}.





\begin{thebibliography}{12}





\bibitem{T0} Тронин С.Н. {\em Операды в категории конвексоров}. I //


Изв. вузов. Математика. -- 2002. -- \No\,3. -- C.\,42--50.





\bibitem{T4} Тронин С.Н. {\em Операды в категории конвексоров} //


Универсальная алгебра и ее приложения. Тез. докл. международн.


семинара, посвящ. памяти проф. Л.А.\,Скорнякова. Волгоград, 6--11 сент.


1999 г. -- Волгоград: ``Перемена'', 1999. -- С.\,62--63.





\bibitem{Max} Максимов В.М. {\em Кубические стохастические матрицы и их


вероятностные интерпретации} // Теория вероятн. и ее примен. -- 1996. --


Т.\,41. -- Bып.\,1. -- С.\,89--106.





\bibitem{T1} Тронин С.Н. {\em О многообразиях, задаваемых полилинейными


тождествами} // Тез. сообщ. XIX Всесоюзн. алгебр. конф. 


9--11 сент. 1987 г. Ч.\,2. Львов, 1987. -- С.\,280.





\bibitem{T2} Тронин С.Н. {\em О некоторых свойствах финитарных алгебраических


теорий} // Тез. сообщ. V Сибирской школы по многообразиям алгебр.


систем, 1--5 июля 1988 г. Барнаул, 1988. -- С.\,68--70.





\bibitem{T3} Тронин С.Н. {\em О некоторых свойствах алгебраических теорий


многообразий линейных ал\-гебр}.~I. {\em Многообразия, задаваемые


полилинейными тождествами}. -- Казанск.


гос. ун-т. -- Казань, 1988. -- 31~с. -- Деп. в ВИНИТИ 11.08.88,


\No\,6511-B88.





\bibitem{T} Тронин С.Н. {\em О ретракциях свободных алгебр и


модулей}: Дис.~$\dotsc$ канд. физ.-матем. наук. Кишинев, 1989. -- 105~с.





\bibitem{Sc1} Скорняков Л.А. {\em Стохастическая алгебра} //


Изв. вузов. Математика. -- 1985. -- \No\,7. -- С.\,3--11.





\bibitem{Sc2} Скорняков Л.А. {\em Алгебра стохастических распределений} //


Изв. вузов. Математика. -- 1982. -- \No\,11. -- С.\,59--67.





\bibitem{Gas} Гаспарян А.С. {\em О некоторых приложениях многомерных


матриц}. -- М.: ВЦ АН СССР, 1983. -- 60~с.





\bibitem{GK} Ginzburg V., Kapranov M. {\em Koszul duality for operads} //


Duke Math. J. -- 1994. -- V.\,76. -- \No\,1. -- P.\,203--272.





\bibitem{Art} Артамонов В.А. {\em Клоны полилинейных операций} //


УMH. -- 1969. -- Т.\,24. -- Bып.\,1. -- С.\,47--59.





\bibitem{May} May J.P. {\em The geometry of iterated loop spaces} //


Lect. Notes Math. -- 1972. -- V.\,271. -- 175~p.





\bibitem{Buh} Бухараев Р.Г. {\em Основы теории вероятностных автоматов}. --


М.: Наука, 1985. -- 400~с.





\bibitem{T5} Тронин С.Н. {\em Многообразия супералгебр и линейные операды} //


Теория функций, ее приложения и смежные вопросы.


Материалы школы-конф., посвящ. 130-летию со дня


рожд. Д.Ф.\,Егорова (Казань, 13--18 сент. 1999 г.) --


Казань: Изд-во Казанск. матем. обществa, 1999. -- С.~224--227.





\end{thebibliography}





\vskip 2cm





\post{\it Казанский государственный университет}{\it Поступила}


\post{}{30.12.1999}





\label{end}


\end{document}








